\newcommand{\tipoRequerimento}{Solicitação de Informações sobre a Taxa Administrativa da CIP}

\section{FUNDAMENTAÇÃO LEGAL DA REQUISIÇÃO}

\begin{enumerate}[resume]
\item
  Para embasar a requisição, utiliza-se as seguintes disposições da Resolução Normativa nº 1.000/2021 (REN nº 1.000/2021) da ANEEL:
  \begin{enumerate}
  \item
    Demonstrativos e Memória de Cálculo (Art. 472, § 1º): A distribuidora tem a obrigação de fornecer ao município os demonstrativos detalhados e a memória de cálculo do faturamento da iluminação pública. Prazos de Atendimento (Art. 408):
    \begin{enumerate}
      \item  
      5 dias úteis para respostas que não exigem visita técnica.
      \item 
      10 dias úteis para os demais casos.
    \end{enumerate}
  \item 
  Prazo Geral (Art. 409): Caso não haja um prazo específico na regulamentação da ANEEL para uma solicitação, a distribuidora deve atendê-la em até 30 dias.
  \end{enumerate}
\end{enumerate}

\section{DOS REQUERIMENTOS}

\begin{enumerate}[resume]
\item Pela atual requisição, solicitamos a concessionária que no prazo máximo de 30 (trinta) dias, encaminhe planilha eletrônica, em formato Excel (.xlsx), contendo as informações seguindo o modelo da tabela, abrangendo o período de \textbf{janeiro de 2020 a dezembro de 2025}; com fundamento no art. 477 da REN nº 1.000/2021:

\begin{table}[ht]
\centering
\caption{Modelo de Planilha Requerida}

{\fontsize{7}{9}\selectfont

\begin{tabularx}{\textwidth}{|
>{\centering\arraybackslash}X|
>{\centering\arraybackslash}X|
>{\centering\arraybackslash}X|
>{\raggedright\arraybackslash}X|
>{\raggedright\arraybackslash}X|
>{\centering\arraybackslash}X|}
\hline

\textbf{PERÍODO DE REFERÊNCIA (MÊS/ANO)} &
\textbf{MONTANTE ARRECADADO DA CIP/COSIP (R\$)} &
\textbf{MONTANTE COBRADO PARA ARRECADAÇÃO DA CIP/COSIP (\%)} &
\textbf{LEGISLAÇÃO QUE CONTÉM PREVISÃO DE ENCONTRO DE CONTAS (NÚMERO DA LEI E DISPOSITIVO)} &
\textbf{MEIO DE COBRANÇA UTILIZADO (FATURA DE ENERGIA, NOTA FISCAL, BOLETO, ETC.)} &
\textbf{REALIZADO ENCONTRO DE CONTAS (SIM OU NÃO)} \\
\hline

... &
... &
... &
... &
... &
... \\
\hline

\end{tabularx}

}

\end{table}

\item Informe se existe legislação municipal vigente que autorize a realização de encontro de contas entre o Município e a distribuidora. Em caso positivo, indicar expressamente o número da lei e o respectivo dispositivo legal (artigo, parágrafo e inciso, se houver).
\item Encaminhe cópia integral da mencionada lei municipal.
\item Esclareça qual a forma atualmente adotada para a cobrança dos valores referentes à arrecadação da CIP/COSIP (por exemplo, inclusão em fatura de energia elétrica, emissão de boleto bancário, nota fiscal ou outro documento de cobrança).
\end{enumerate}